\section{Introduction}
\label{sec:intro}

Successful long-horizon manipulation requires connecting what has already happened to what should happen next.
World--action models (WAMs) couple visual prediction with action generation, bringing predictive representations of physical interaction into robot control~\citep{kim2026cosmospolicy,li2026lingbotva,bi2025motus}.
Yet understanding how a scene can change does not determine which change the current task calls for, or which past behavior is worth reusing.

Existing approaches address different parts of this connection.
\emph{Local-context WAMs} act from current observations or short histories (Figure~\ref{fig:leadin}a), which can omit earlier interactions needed to disambiguate the next action~\citep{kaelbling1998pomdp,dai2026robomme,lei2026robomemarena}.
\emph{Memory-augmented WAMs} retain historical context or derive memory-grounded plans (Figure~\ref{fig:leadin}b)~\citep{yang2026memorywam,wang2026dimwam,zhao2026mapwam}.
Recent VLA policies also combine online visual history with retrieved experience prompts~\citep{zhou2026remembersmarter}.
Access to either alone does not establish a recalled action's present applicability.
\emph{Trajectory retrieval and informed flow sources} supply action exemplars and structured generation priors~\citep{park2026recap,kang2026warmprior,jia2026a2a,machado2026lafm}; a historically successful action, however, may repeat a completed interaction or produce an unwanted change.
These approaches raise a common question: \textbf{\emph{how can a WAM infer the current task requirement, identify applicable past actions, and regulate their influence on generation?}}


Our perspective echoes the \emph{Analects}' idea of \textit{gaining new insight by revisiting prior knowledge}~\citep[2.11]{confucius2002analects}.
\emph{Recollection} retains observed changes and executed-action summaries from the current episode; a \emph{repertoire} preserves past behaviors and their observed effects across episodes.
The former informs current task requirements; the latter offers reusable action structure rather than merely additional context.
Familiar behavior alone cannot reveal what has already happened here, so its applicability must be reassessed against current evidence.
Our insight is to use \emph{recollection as evidence for assessing the repertoire, and recalled behavior as a revisable action prior}.

\begin{figure}[!t]
  \centering
  \includegraphics[width=1.00\linewidth]{fig/overview.pdf}
  \caption{\textbf{From interaction evidence to selective action reuse.}
(a,b) Illustrative WAM interfaces use local observations or retained history.
(c) Recollection grounds the task requirement; the repertoire supplies action--effect proposals.
Predicted consequence compatibility informs selection and a shared gate over source and conditioning.
K/V denotes visual keys and values; dashed arrows indicate training-only prediction.
\textbf{Right:} Benchmark results with published references (Tables~\ref{tab:libero}--\ref{tab:rmbench}) and distraction on Put Back Block and Press Button, averaged equally within each of five evaluation runs.}
  \label{fig:leadin}
\vspace{-0.3cm}
\end{figure}
\vspace{-0.3cm}

We introduce \method{}, the \emph{World-Action-Retrospection Model}, a video--action architecture for \emph{event-grounded retrospection} (Figure~\ref{fig:leadin}c).
Current world features and \emph{situated recollection} ground the task requirement; the \emph{sensorimotor repertoire} supplies cross-episode state--action--effect events as reusable proposals.
We adapt and assess these proposals by predicted effect, then use \emph{consequence-gated source transport} to regulate their contribution to Action DiT initialization and conditioning.

\emph{Consequence-aligned action recall} compares the required change with each recalled action's predicted effect.
We infer a candidate-independent requirement from current context and recollection, and predict proposal effects from adapted actions, world features, and stored effect evidence.
Compatibility complements history-aware ranking and guides selection and gating.

\emph{Consequence-gated source transport} carries this assessment into generation.
A shared gate jointly regulates proposal-informed stochastic initialization and candidate-derived conditioning, letting Action DiT refine useful action structure as a revisable prior.
Closing the gate restores Gaussian initialization while preserving recollection-grounded task conditioning.
This shared control admits, attenuates, or withdraws recalled behavior without iterative future-video generation.

Our contributions are threefold:
\begin{itemize}
    \item We introduce event-grounded retrospection for WAMs, linking history-dependent task understanding to selective reuse of cross-episode action experience.
    \item We develop consequence-aligned action recall and consequence-gated source transport: a candidate-independent requirement guides proposal assessment, while a shared gate controls initialization and conditioning for selective generative refinement.
    \item \method{} achieves 98.7\% success on LIBERO~\citep{liu2023libero} and 83.9\% on RMBench~\citep{chen2026rmbench}, including 87.3\% on $M(n)$ tasks requiring multiple past observations.
    Component results suggest complementary benefits of recollection and repertoire; complete-event distraction tests support explicit comparison with partial event replacement.
\end{itemize}
